\subsection{Max-rank of matrix subspaces and slice spans of a tensor}\label{subsec:maxrank}

%
    %
    We denote by $\FF^{n_1} \otimes \FF^{n_2}$ the space of $n_1 \times n_2$ matrices over $\FF$.
    We define the \emph{max-rank} of any matrix subspace $\mathcal{A} \subseteq \FF^{n_1} \otimes \FF^{n_2}$ as
    %
    \[
\maxrank(\mathcal{A}):=\max\{\rank(A):A\in\mathcal{A}\}.
    \]
    %
    %
    %
    %
%
%
For any tensor $T = \sum_{i,j,k} T_{i,j,k}\, e_i \otimes e_j \otimes e_k \in \FF^{n_1} \otimes \FF^{n_2} \otimes \FF^{n_3}$ we define the matrices %
%
%
%
%
%
%
\begin{align*}
T^{(1)}_i &= \sum_{j,k} T_{i,j,k}\, e_j \otimes e_k \quad (i \in [n_1])\\
T^{(2)}_j &= \sum_{i,k} T_{i,j,k}\, e_i \otimes e_k \quad (j \in [n_2])\\
T^{(3)}_k &= \sum_{i,j} T_{i,j,k}\, e_i \otimes e_j \quad (k \in [n_3]).
\end{align*}
For every $\ell \in [3]$, we call $\{T^{(\ell)}_i : i \in [n_\ell] \}$ the \emph{$\ell$-slices} of $T$. %
%
%
%
%
%
%
%
%
%
%
%
%
%
%
%
%
%
%
%
%
%
%
%
%
%
%
%
%
%
%
%
%
%
%
%
%
%
%
%
%
%
%
We define the $\ell$-slice span of $T$ as the matrix subspace
$\mathcal{A}_\ell = \linspan \{T^{(\ell)}_i : i \in [n_\ell] \}$
and denote its max-rank by
\[
\subrank_\ell(T) = \maxrank(\mathcal{A}_\ell).
\]
%
%
We note that the flattening rank $\tensorrank_\ell(T)$ defined earlier, is equal to the dimension of the subspace $\mathcal{A}_\ell$. 

%



%
%



%
%

%
%
%
%
%
%
%
%

%
%
%

%
%
%
%
%
%
%
%
%
%
%
%
%
%
%
%

%

%
%
%
%
%
%
%
%
%
%
%
%
%




\subsection{Lower bound on products of max-ranks of slice spans}\label{subsec:QiQj}

%

%

%

%

%\refnote{9}{The following paragraph sounds like a repetitions of earlier things. also \phantomsection\label{ref:9-repetition-22}\ref{ref:9-repetition-21}}
In this section we prove the aforementioned fundamental property of the max-ranks of the slice spans of a tensor, showing that they cannot all be small. This will be a crucial ingredient for proving the asymptotic subrank lower bound afterwards.

\begin{theorem}\label{th:uncertainty-principle-n}
    Let $T \in \FF^{n_1} \otimes \FF^{n_2} \otimes \FF^{n_3}$ be concise. Let ${\ell_1}, {\ell_2}, {\ell_3} \in [3]$ be distinct. Suppose that $|\FF|>n_{\ell_3}$. Then 
    \[
    \subrank_{\ell_1}(T)\subrank_{\ell_2}(T) \geq n_{\ell_3}.
    \]
\end{theorem}

%

%

%

A corollary of this is that there are two distinct directions with large max-rank:

\begin{corollary} %
\label{cor:twoLargeRankDirs}
    Let $T\in\FF^{n_1}\otimes\FF^{n_2}\otimes\FF^{n_3}$ be concise.
    Suppose $n_1 \leq n_2 \leq n_3$. Suppose that $|\FF| > n_3$. 
    %
    Then there are distinct ${\ell_1},{\ell_2}\in[3]$ such that
    \begin{align*}\subrank_{\ell_1}(T) &\geq \sqrt{n_3},\\
    \subrank_{\ell_2}(T) &\geq \sqrt{n_1}.
    \end{align*}
\end{corollary}
%
\begin{proof}%
    %
    We have by \autoref{th:uncertainty-principle-n} that
    \begin{align*}
    \subrank_1(T) \subrank_2(T) &\geq n_3\\
    \subrank_2(T) \subrank_3(T) &\geq n_1\\
    \subrank_1(T) \subrank_3(T) &\geq n_2 \geq n_1.
    \end{align*}
    Thus one of $\subrank_1(T)$ and $\subrank_2(T)$ is at least $\sqrt{n_3}$, and either $\subrank_3(T)$ is at least~$\sqrt{n_1}$ or both $\subrank_1(T)$ and $\subrank_2(T)$ are at least $\sqrt{n_1}$.
\end{proof}
